% !TEX program = platex
% 数学ノート集・共通テンプレート v2（積分論2026年度第1部の版面を参考）
% A5判。章内で定義・定理・例・注意等の番号を共通化。
% コンパイル： platex -interaction=nonstopmode math_note_template_v2.tex (2回)
%              dvipdfmx math_note_template_v2.dvi
\documentclass[a5paper,10pt,onecolumn]{jsarticle}
\usepackage[dvipdfmx,top=22mm,bottom=20mm,inner=19mm,outer=18mm,headheight=13pt,headsep=6mm]{geometry}
\usepackage{amsmath,amssymb,amsthm,mathtools}
\usepackage{enumitem,fancyhdr,titlesec}
\usepackage[dvipdfmx,bookmarks=true,colorlinks=false]{hyperref}
\usepackage{pxjahyper}

% ----- 差し替える文書情報 -----
\newcommand{\NoteTitle}{Lipschitz領域のSobolevトレース定理}
\newcommand{\NoteSeries}{数学ノート集}

% ----- 文字・組版 -----
% 参考PDFの質素な明朝＋Computer Modernの印象を引き継ぎながら，
% A5判で窮屈にならないよう本文の行間を少し確保する．
\linespread{1.13}
\setlength{\parindent}{1zw}
\setlength{\parskip}{0pt}
\setlength{\abovedisplayskip}{6pt plus 2pt minus 2pt}
\setlength{\belowdisplayskip}{6pt plus 2pt minus 2pt}
\setlength{\abovedisplayshortskip}{3pt plus 2pt minus 1pt}
\setlength{\belowdisplayshortskip}{4pt plus 2pt minus 1pt}
\setlist[enumerate]{topsep=3pt,itemsep=1pt,parsep=0pt,leftmargin=2.4zw}
\titleformat{\section}{\normalfont\normalsize\bfseries}{\thesection}{0.8zw}{}
\titlespacing*{\section}{0pt}{2.0ex plus .4ex}{.9ex}
\titleformat{\subsection}{\normalfont\normalsize\bfseries}{\thesubsection}{0.7zw}{}
\titlespacing*{\subsection}{0pt}{1.5ex plus .3ex}{.6ex}
\raggedbottom

% ----- 柱とページ番号：ページ上部、罫線なし -----
\fancypagestyle{mathnote}{%
  \fancyhf{}%
  \fancyhead[L]{\footnotesize \NoteSeries\ （\NoteTitle）}%
  \fancyhead[R]{\footnotesize \thepage}%
  \renewcommand{\headrulewidth}{0pt}%
  \renewcommand{\footrulewidth}{0pt}%
}
\pagestyle{mathnote}

% ----- 節内通し番号の定義・定理群 -----
\newtheoremstyle{mathnotetheorem}
 {4.5pt}{4.5pt}{\normalfont}{}%
 {\bfseries}{}{0.65zw}%
 {\thmname{#1}\thmnumber{ #2}\thmnote{（#3）}}
\theoremstyle{mathnotetheorem}
\newtheorem{theorem}{定理}[section]
\newtheorem{definition}[theorem]{定義}
\newtheorem{lemma}[theorem]{補題}
\newtheorem{proposition}[theorem]{命題}
\newtheorem{corollary}[theorem]{系}
\newtheorem{example}[theorem]{例}
\newtheorem{remark}[theorem]{注意}
\renewcommand{\proofname}{証明}
\renewcommand{\qedsymbol}{$\square$}
\makeatletter
\renewenvironment{proof}[1][\proofname]{%
  \par\pushQED{\qed}\normalfont\topsep4\p@\@plus2\p@\relax
  \trivlist\item[\hskip\labelsep\bfseries #1]\ignorespaces
}{\popQED\endtrivlist\@endpefalse}
\makeatother

% ----- 使うノートで適宜拡張する記号 -----
\newcommand{\R}{\mathbb{R}}
\newcommand{\norm}[1]{\lVert #1\rVert}
\newcommand{\abs}[1]{\lvert #1\rvert}

\begin{document}
\section*{読了後にできること}
本稿を読むための前提は，Sobolev空間の弱微分，Fourier変換のPlancherel定理，
標準mollifier，Lipschitz写像の変数変換である．トレース作用素自体は本稿で構成する．
本稿の具体的な到達課題は，一次元のトレースを具体的に計算し，零トレースの意味を確認する．

\section{問題の定式化}
本稿では，有界Lipschitz領域 $\Omega\subset\mathbb R^n$（$n\ge2$）における
Sobolevトレース定理を証明する．
主結果は連続な全射
\[
 \operatorname{Tr}:H^1(\Omega)\longrightarrow H^{1/2}(\partial\Omega)
\]
と，その有界な右逆作用素の構成，および
\[
 \ker\operatorname{Tr}=H_0^1(\Omega)
\]
である．必要な局所延長と稠密性を，境界の平坦化を用いて示す．

前提にする基本事項は，Sobolev弱微分，Fourier変換のPlancherel定理，
標準mollifierによる近似，Lipschitz写像のほとんど至る所での微分可能性
（Rademacherの定理），変数変換公式，および有限分割単位の存在である．
これらは境界トレース定理そのものを前提とするものではない．

\begin{definition}[Lipschitz領域]
有界開集合 $\Omega\subset\mathbb R^n$ がLipschitz領域であるとは，
任意の $x_0\in\partial\Omega$ の近傍で座標を剛体運動により選び直すと，
ある $a>0$，開球 $B\subset\mathbb R^{n-1}$，Lipschitz函数
$\varphi:B\to\mathbb R$ が存在して，局所的に
\[
 \Omega=\{(x',t):x'\in B,\ t>\varphi(x')\}
\]
と表されることをいう．有限個のこのような境界近傍で
$\partial\Omega$ を覆える．
\end{definition}

\begin{remark}
各境界座標は，境界までの距離の比較に必要な余裕を持って
少し大きい座標近傍を用意しておく．
以下の局所的な函数は，この余裕の内側に台を置く．
この処理により，垂直方向の平行移動や畳み込みを
近傍の側面境界に干渉せずに行える．
\end{remark}

\section{半空間の一次元評価とFourierトレース}
\begin{lemma}[半直線の端点値]
$v\in H^1(0,\infty)$ は連続な代表を持ち，任意の $a>0$ に対して
\[
 a|v(0)|^2\le\int_0^\infty(a^2|v(t)|^2+|v'(t)|^2)\,dt
\]
が成り立つ．
\end{lemma}
\begin{proof}
一次元Sobolev函数はコンパクト区間上絶対連続な代表を持つ．
$v\in L^2(0,\infty)$ だから $t_k\to\infty$ で
$v(t_k)\to0$ となる部分列が取れる．
$|v|^2$ の積分公式から
\[
 |v(0)|^2=-2\operatorname{Re}\int_0^\infty
     v'(t)\overline{v(t)}\,dt
\]
を得る．右辺はCauchy--Schwarzで絶対収束する．
$2a|v'||v|\le|v'|^2+a^2|v|^2$ を積分すればよい．
\end{proof}

$\mathbb R^n_+=\mathbb R^{n-1}\times(0,\infty)$，$m=n-1$ と置く．
$H^{1/2}(\mathbb R^m)$ は
\[
 \norm g_{H^{1/2}(\mathbb R^m)}^2
 =\int_{\mathbb R^m}(1+|\xi|^2)^{1/2}|\widehat g(\xi)|^2\,d\xi
\]
をノルムとするHilbert空間として定義する．

\begin{lemma}[半空間上の滑らかな函数の稠密性]
\label{lem:halfspace-density}
$H^1(\mathbb R^n_+)$ において
$C_c^\infty(\overline{\mathbb R^n_+})$ は稠密である．
ここで右辺は $\mathbb R^n$ 上の滑らかな
コンパクト台函数の半空間への制限全体を表す．
\end{lemma}
\begin{proof}
まず $u\in H^1(\mathbb R^n_+)$ とする．
$\chi\in C_c^\infty(\mathbb R^n)$ を
$0\le\chi\le1$，
$\chi=1$ on $B(0,1)$，
$\operatorname{supp}\chi\subset B(0,2)$
となるように取り，$\chi_R(x)=\chi(x/R)$ と置く．
積の弱微分公式から
\[
 \partial_j(\chi_Ru)
 =\chi_R\partial_ju+(\partial_j\chi_R)u .
\]
優収束定理と
$\norm{\nabla\chi_R}_\infty\le C/R$ により
$\chi_Ru\to u$ が $H^1(\mathbb R^n_+)$ で成り立つ．
よって半空間の閉包内で有界な台を持つ函数を
近似すれば十分である．

そのような函数を $v$ とする．
$v$ および各弱導函数を $L^2$ 函数として
全空間に零延長したものを，それぞれ
$v^0,(\partial_jv)^0$ と書く．
ここでは $v^0$ が全空間の $H^1$ に属するとは
仮定しない．
$\varepsilon>0$ とし
\[
 v_\varepsilon(x',t)=v(x',t+\varepsilon)
 \qquad(t>-\varepsilon)
\]
と定める．これは
$\mathbb R^{n-1}\times(-\varepsilon,\infty)$ 上の
$H^1$ 函数であり，弱導函数は対応する
$\partial_jv$ の平行移動に等しい．

$\theta\in C^\infty(\mathbb R)$ を
\[
 \theta(t)=0\ (t\le-3/4),\qquad
 \theta(t)=1\ (t\ge-1/2)
\]
となるように取る．
$t>-\varepsilon$ で
$W_\varepsilon(x',t)=\theta(t/\varepsilon)
 v_\varepsilon(x',t)$ と置き，それ以外では零とする．
切断函数の積の弱微分公式から，
$W_\varepsilon\in H^1(\mathbb R^n)$ である．
実際，$W_\varepsilon$ は
$t\le-3\varepsilon/4$ で恒等的に零であるため，
延長の境界に跳びは生じない．
一方 $t>0$ では $\theta(t/\varepsilon)=1$ だから，
\[
 W_\varepsilon(x',t)=v(x',t+\varepsilon),\qquad
 \partial_jW_\varepsilon(x',t)
 =(\partial_jv)(x',t+\varepsilon)
\]
がほとんど至る所成り立つ．
$L^2(\mathbb R^n)$ の平行移動連続性を
$v^0$ と $(\partial_jv)^0$ に適用すると
\[
 \norm{W_\varepsilon|_{\mathbb R^n_+}-v}
       _{H^1(\mathbb R^n_+)}
 \longrightarrow0
 \qquad(\varepsilon\downarrow0)
\]
を得る．

$\varepsilon$ を固定して標準mollifierを取れば，
$\rho_\delta*W_\varepsilon$ は
$\mathbb R^n$ 上の滑らかなコンパクト台函数であり，
$\delta\downarrow0$ で $W_\varepsilon$ に
$H^1(\mathbb R^n)$ 収束する．
その半空間への制限を取ってから，
$\varepsilon$ と $\delta$ を順に小さくすれば，
$v$ を $C_c^\infty(\overline{\mathbb R^n_+})$ で
$H^1$ 近似できる．最後に $R\to\infty$ とする．
\end{proof}

\begin{theorem}[半空間のトレース不等式]
$u\in C_c^\infty(\overline{\mathbb R^n_+})$ とし，
$g(x')=u(x',0)$ とする．このとき
\[
 \norm g_{H^{1/2}(\mathbb R^m)}
 \le\norm u_{H^1(\mathbb R^n_+)}.
\]
従って境界への制限は，連続作用素
$\operatorname{Tr}_+:H^1(\mathbb R^n_+)\to H^{1/2}(\mathbb R^m)$
へ一意に延長される．
\end{theorem}
\begin{proof}
$x'$ に関するFourier変換を $\widehat u(\xi,t)$ とする．
前補題を各 $\xi$ について
$a=(1+|\xi|^2)^{1/2}$ で適用すると
\[
 (1+|\xi|^2)^{1/2}|\widehat u(\xi,0)|^2
 \le\int_0^\infty
 \{(1+|\xi|^2)|\widehat u(\xi,t)|^2
                 +|\partial_t\widehat u(\xi,t)|^2\}\,dt.
\]
Plancherelの等式とTonelliの定理によって $\xi$ で積分すればよい．
補題\ref{lem:halfspace-density}により
$C_c^\infty(\overline{\mathbb R^n_+})$ は
$H^1(\mathbb R^n_+)$ に稠密である．
従って不等式はこの近似列に沿って極限へ移り，
境界制限の一意な有界線形延長が得られる．
\end{proof}

\begin{lemma}[半空間の境界値と平行断面]
\label{lem:halfspace-slice}
任意の $w\in H^1(\mathbb R^n_+)$ は，
$t\ge0$ に関して $L^2(\mathbb R^{n-1})$ 値連続な
代表 $t\mapsto w_t$ を持つ．さらに
\[
 w_t\longrightarrow\operatorname{Tr}_+w
 \quad\text{in }L^2(\mathbb R^{n-1})
 \qquad(t\downarrow0)
\]
が成り立つ．
\end{lemma}
\begin{proof}
まず $w\in C_c^\infty(\overline{\mathbb R^n_+})$ とする．
任意の $s,t\ge0$ に対して
\[
 \norm{w_t-w_s}_{L^2}
 \le |t-s|^{1/2}\norm{\partial_tw}_{L^2(\mathbb R^n_+)}
\]
が微積分学の基本定理とCauchy--Schwarzから従う．
また半直線の端点値の補題を
$v(r)=w(x',t+r)$，$a=1$ に適用して
$x'$ について積分すると
\[
 \sup_{t\ge0}\norm{w_t}_{L^2}^2
 \le\norm{w}_{H^1(\mathbb R^n_+)}^2
\]
を得る．
一般の $w$ に対して
補題\ref{lem:halfspace-density}により
$w^{(k)}\in C_c^\infty(\overline{\mathbb R^n_+})$，
$w^{(k)}\to w$ in $H^1$ を取る．
上の上限評価を差に適用すると，
$w^{(k)}_t$ は $t\ge0$ にわたり
$L^2$ で一様Cauchyとなる．
従って $L^2$ 値連続な極限 $w_t$ を持ち，
その代表はほとんど全ての $t>0$ で元の $w$ に一致する．
トレース不等式と $H^1$ 収束から
$w^{(k)}_0\to\operatorname{Tr}_+w$ in $L^2$ である．
一方，断面の一様収束から
$w^{(k)}_0\to w_0$ in $L^2$ なので，
$w_0=\operatorname{Tr}_+w$ を得る．
$t\mapsto w_t$ の連続性により結論が従う．
\end{proof}

\begin{theorem}[半空間のトレースの右逆]
$g\in H^{1/2}(\mathbb R^m)$ に対し，
$x'$ に関するFourier変換を用いて
\[
 \widehat{Eg}(\xi,t)=e^{-t(1+|\xi|^2)^{1/2}}\widehat g(\xi)
 \qquad(t>0)
\]
と定める．このとき $Eg\in H^1(\mathbb R^n_+)$ であり，
$\operatorname{Tr}_+(Eg)=g$，
\[
 \norm{Eg}_{H^1(\mathbb R^n_+)}^2
 =\norm g_{H^{1/2}(\mathbb R^m)}^2
\]
が成立する．
\end{theorem}
\begin{proof}
$a=(1+|\xi|^2)^{1/2}$ とする．
PlancherelとTonelliの定理によって
\[
 \begin{aligned}
 \norm{Eg}_{H^1}^2
 &=\int_{\mathbb R^m}\int_0^\infty
   (a^2+a^2)e^{-2at}|\widehat g(\xi)|^2\,dt\,d\xi\\
 &=\int_{\mathbb R^m}a|\widehat g(\xi)|^2\,d\xi.
 \end{aligned}
\]
さらに $Eg$ の $L^2$ 値断面について
\[
 \norm{(Eg)_t-g}_{L^2}^2
 =\int_{\mathbb R^m}
   |e^{-t(1+|\xi|^2)^{1/2}}-1|^2
   |\widehat g(\xi)|^2\,d\xi
 \longrightarrow0\qquad(t\downarrow0)
\]
となる．実際，被積分函数は
$4|\widehat g(\xi)|^2$ 以下であり，
優収束定理が使える．
補題\ref{lem:halfspace-slice}により同じ極限は
$\operatorname{Tr}_+(Eg)$ に等しいから，
$\operatorname{Tr}_+(Eg)=g$ を得る．
\end{proof}

\section{分数階の差分による表示}
$H^{1/2}(\mathbb R^m)$ のノルムは局所座標の変換に都合のよい
差分積分によっても表される．

\begin{lemma}[Gagliardo半ノルムとFourier乗数]
$g\in H^{1/2}(\mathbb R^m)$ に対し
\[
 [g]_{1/2,\mathbb R^m}^2
 =\iint_{\mathbb R^m\times\mathbb R^m}
    \frac{|g(x)-g(y)|^2}{|x-y|^{m+1}}\,dx\,dy
\]
と定めると，ある $c_m>0$ が存在して
\[
 [g]_{1/2,\mathbb R^m}^2
 =c_m\int_{\mathbb R^m}|\xi|\,|\widehat g(\xi)|^2\,d\xi.
\]
従って $\norm g_{L^2}^2+[g]_{1/2}^2$ は
$H^{1/2}$ ノルムと同値である．
\end{lemma}
\begin{proof}
まずSchwartz函数 $g$ を取る．$h=x-y$ と変数変換し，
Plancherelを用いると
\[
 [g]_{1/2}^2
 =\int_{\mathbb R^m}|\widehat g(\xi)|^2
  \left(\int_{\mathbb R^m}
      \frac{|e^{i\xi\cdot h}-1|^2}{|h|^{m+1}}\,dh\right)d\xi.
\]
内側の積分は回転不変性と $h\mapsto|\xi|^{-1}h$ により
$c_m|\xi|$ に等しい．$c_m$ が正で有限なことは
$|e^{it}-1|\le\min(2,|t|)$ を使えばよい．
実際，原点近傍の動径積分は有界で，無限遠の動径積分は
$\int_1^\infty r^{-2}dr<\infty$ となる．
一般の $g\in H^{1/2}(\mathbb R^m)$ を取る．
Schwartz函数は $H^{1/2}$ で稠密である：
実際，Fourier側の重み付き $L^2$ 空間で
コンパクト台の滑らかな函数を近似し，
Fourier逆変換を取ればよい．
従って $g_j\in\mathcal S$ で
$g_j\to g$ in $H^{1/2}$ となる列を取れる．
Schwartz函数に対して証明した等式を $g_j-g_k$ に適用すると
\[
 \iint
 \frac{|(g_j-g_k)(x)-(g_j-g_k)(y)|^2}{|x-y|^{m+1}}
 \,dx\,dy
 =c_m\int|\xi|\,|\widehat g_j-\widehat g_k|^2\,d\xi
 \longrightarrow0.
\]
よって
$D_j(x,y)=(g_j(x)-g_j(y))/|x-y|^{(m+1)/2}$
は $L^2(\mathbb R^m\times\mathbb R^m)$ のCauchy列である．
その極限を $D$ とする．
一方 $g_j\to g$ in $L^2$ なので，
部分列を取れば $g_j(x)\to g(x)$ がほとんど至る所で成立する．
さらに必要なら部分列を取り直すと
$D_j(x,y)\to D(x,y)$ がほとんど至る所で成立する．
従って対角集合を除くほとんど全ての $(x,y)$ について
\[
 D(x,y)=\frac{g(x)-g(y)}{|x-y|^{(m+1)/2}}
\]
である．
よって $g$ のGagliardo半ノルムは有限で，
$D_j\to D$ の $L^2$ 収束と
$\widehat g_j\to\widehat g$ の重み付き $L^2$ 収束から
求める等式を得る．
最後に $1+|\xi|$ と $(1+|\xi|^2)^{1/2}$ の
二側評価からノルムの同値性が従う．
\end{proof}

\section{Lipschitz座標でのSobolevノルム}
境界を平坦化する写像を
\[
 \Phi(x',t)=(x',t+\varphi(x'))
\]
とする．ここで $\varphi$ はLipschitz函数である．
$\Phi$ と $\Phi^{-1}$ はともにLipschitzであり，
$\det D\Phi=1$ がほとんど至る所で成立する．

\begin{lemma}[局所の座標変換]
\label{lem:chart-H1}
$B\subset\mathbb R^{n-1}$ を球，$h>0$ とし，
$\varphi:B\to\mathbb R$ をLipschitz函数とする．
\[
 Q_\varphi=\Phi(B\times(-h,h)),\qquad
 \Phi(x',s)=(x',s+\varphi(x'))
\]
が境界座標近傍であり，
$\Omega\cap Q_\varphi=\Phi(B\times(0,h))$
であるとする．
$\chi\in C_c^\infty(Q_\varphi)$ を取る．
$u\in H^1(\Omega)$ に対して
\[
 v(x',s)=(\chi u)(\Phi(x',s))
 \quad ((x',s)\in B\times(0,h))
\]
を定め，それ以外の $\mathbb R^n_+$ で零とすると，
$v\in H^1(\mathbb R^n_+)$ で
\[
 \norm v_{H^1(\mathbb R^n_+)}
 \le C_\chi\norm u_{H^1(\Omega)}
\]
が成り立つ．
逆に $w\in H^1(\mathbb R^n_+)$ について
\[
 Uw(x)=
 \begin{cases}
 \chi(x)w(\Phi^{-1}(x))&x\in\Omega\cap Q_\varphi,\\
 0&x\in\Omega\setminus Q_\varphi
 \end{cases}
\]
とすれば $Uw\in H^1(\Omega)$ かつ
$\norm{Uw}_{H^1(\Omega)}
 \le C_\chi\norm w_{H^1(\mathbb R^n_+)}$ である．
\end{lemma}
\begin{proof}
$\Phi$ と $\Phi^{-1}$ はbi-Lipschitzであり，
Rademacherの定理によりほとんど至る所で
\[
 D\Phi=
 \begin{pmatrix}I_{n-1}&0\\ \nabla\varphi&1\end{pmatrix},
 \qquad
 \det D\Phi=1
\]
が成り立つ．その微分と逆写像の微分は
Lipschitz定数に応じて一様有界である．

まず $u$ が $\overline\Omega$ の近傍で滑らかなとき，
$\chi u$ は境界座標内で滑らかかつ有界な導函数を持つ．
Lipschitz写像との合成に対する連鎖律より
\[
 \nabla((\chi u)\circ\Phi)
 =D\Phi^{\,T}(\nabla(\chi u))\circ\Phi
 \quad\text{a.e.}
\]
である．変数変換公式と
$|\nabla(\chi u)|\le
 |\chi||\nabla u|+|\nabla\chi||u|$
から上記のノルム評価を得る．
$\operatorname{supp}\chi\Subset Q_\varphi$ なので，
平坦座標の側面 $x'\in\partial B$ と上面 $s=h$ の
近傍で $v$ は零になる．従って人工的な境界を
横断する零延長は $H^1$ に属する．
一般の $u\in H^1(\Omega)$ は
後掲のLipschitz領域での稠密性補題
\ref{lem:lip-density}で滑らかな函数に近似する．
ノルム評価から平坦化した列が $H^1$ で収束し，
その極限は $L^2$ の同定から上記の $v$ に等しい．

逆向きには，まず
$w\in C_c^\infty(\overline{\mathbb R^n_+})$ を取る．
$w\circ\Phi^{-1}$ は局所的にLipschitzであり，
連鎖律と変数変換によって
$\norm{Uw}_{H^1(\Omega)}
 \le C_\chi\norm w_{H^1(\mathbb R^n_+)}$
が従う．
$\chi$ の台は $Q_\varphi$ の内部にあるから，
$Q_\varphi$ の外への零延長で
人工的な境界項は生じない．
一般の $w$ については
補題\ref{lem:halfspace-density}による近似と
上記の評価で極限を取り，
$L^2$ で同定すればよい．
\end{proof}

\begin{lemma}[Lipschitz領域上の滑らかな函数の稠密性]
\label{lem:lip-density}
有界Lipschitz領域 $\Omega$ に対して
$C^\infty(\overline\Omega)$ は $H^1(\Omega)$ で稠密である．
\end{lemma}
\begin{proof}
境界の有限個のグラフ座標近傍と内部の開集合を取り，
それらに従属する滑らかな分割単位
$\{\eta_j\}_{j=0}^N$ を選ぶ．
$u_j=\eta_j u$ と置けば $u=\sum_j u_j$ であり，
各 $u_j$ の台は対応する座標近傍の内部に
コンパクトに含まれるように選べる．
内部に台を持つ成分は零延長と通常のmollifierで
$H^1$ 近似できる．以下，境界成分一つを考える．

直交座標 $(x',t)$ で局所的に
$\Omega=\{t>\varphi(x')\}$ と表し，
$u_j$ の台と座標近傍の側面との間には
正の距離があるものとする．十分小さい
$\varepsilon>0$ について，
\[
 v_\varepsilon(x',t)=u_j(x',t+\varepsilon)
 \qquad(t>\varphi(x')-\varepsilon)
\]
と定める．平行移動のもとで弱微分は
$\partial_kv_\varepsilon
 =(\partial_ku_j)(x',t+\varepsilon)$
となるので，$v_\varepsilon$ は表示された領域の
$H^1$ 函数である．
$\theta\in C^\infty(\mathbb R)$ を
\[
 \theta(s)=0\quad(s\le-3/4),\qquad
 \theta(s)=1\quad(s\ge-1/2)
\]
となるように取る．
$t>\varphi(x')-\varepsilon$ で
\[
 W_\varepsilon(x',t)
 =\theta\!\left(\frac{t-\varphi(x')}{\varepsilon}\right)
  v_\varepsilon(x',t)
\]
と置き，それ以外では零と定める．
$\varphi$ はLipschitz函数なので，
$\theta((t-\varphi(x'))/\varepsilon)$ は
$W^{1,\infty}$ 函数である．
しかも $W_\varepsilon$ は
$t\le\varphi(x')-3\varepsilon/4$ で零であり，
台は座標近傍の内部に留まる．
従って積の弱微分公式と零延長によって
$W_\varepsilon\in H^1(\mathbb R^n)$ が得られる．
元の領域側では切断函数が恒等的に1なので
\[
 W_\varepsilon|_\Omega
 =u_j(x',t+\varepsilon),\qquad
 \nabla W_\varepsilon|_\Omega
 =(\nabla u_j)(x',t+\varepsilon)
 \quad\text{a.e.}
\]
である．

ここで $u_j$ とその各弱導函数を，
単に $L^2$ 函数として $\mathbb R^n$ 上に零延長する．
この零延長の弱微分が元の弱微分の零延長に等しいとは
主張していない．$L^2(\mathbb R^n)$ における
平行移動の連続性だけを用いれば，
上の領域内部での微分表示から
\[
 \norm{W_\varepsilon|_\Omega-u_j}_{H^1(\Omega)}
 \longrightarrow0\qquad(\varepsilon\downarrow0)
\]
が従う．

$\varepsilon$ を固定すると
$W_\varepsilon\in H^1(\mathbb R^n)$ だから，
標準mollifierにより
$\rho_\delta*W_\varepsilon\to W_\varepsilon$
が $\delta\downarrow0$ で $H^1(\mathbb R^n)$ の意味で成り立つ．
$\delta$ を座標近傍の台の余裕より小さく選べば，
$\rho_\delta*W_\varepsilon$ はその近傍に台を持つ
$C_c^\infty(\mathbb R^n)$ 函数である．
従って三角不等式から
\[
 \norm{(\rho_\delta*W_\varepsilon)|_\Omega-u_j}_{H^1(\Omega)}
 \le
 \norm{\rho_\delta*W_\varepsilon-W_\varepsilon}_{H^1(\mathbb R^n)}
 +\norm{W_\varepsilon|_\Omega-u_j}_{H^1(\Omega)}
\]
となる．最初に $\varepsilon$ を，次に $\delta$ を選び，
右辺を任意に小さくできる．
有限個の成分について得た滑らかな近似を足せば，
$u$ は $\overline\Omega$ の近傍で滑らかな函数によって
$H^1(\Omega)$ 近似される．
\end{proof}

\section{境界の分数階Sobolev空間}
$\partial\Omega$ には各Lipschitz座標の
$g\mapsto g\circ(x'\mapsto(x',\varphi(x')))$ という対応がある．
境界面積測度を $dS$ とする．

\begin{definition}
$H^{1/2}(\partial\Omega)$ を $g\in L^2(\partial\Omega,dS)$ で
\[
 \norm g_{H^{1/2}(\partial\Omega)}^2
 =\norm g_{L^2(\partial\Omega)}^2
 +\iint_{\partial\Omega\times\partial\Omega}
   \frac{|g(x)-g(y)|^2}{|x-y|^n}\,dS_x\,dS_y<\infty
\]
を満たすものの空間とする．
\end{definition}

\begin{lemma}[境界分数階ノルムの局所化]
\label{lem:boundary-local}
$\eta$ を $\partial\Omega$ 上のLipschitz函数とする．
このとき
\[
 \norm{\eta g}_{H^{1/2}(\partial\Omega)}
 \le C_\eta\norm g_{H^{1/2}(\partial\Omega)}.
\]
さらに，$V\subset\partial\Omega$ を
グラフ座標 $\psi:B\to V$，
$\psi(x')=(x',\varphi(x'))$ で表される境界部分とし，
$K\Subset B$ をコンパクト集合とする．
$\operatorname{supp}(\eta|_{\partial\Omega})\subset\psi(K)$
なら，
\[
 T_\eta g(x')=
 \begin{cases}
  (\eta g)(\psi(x'))&x'\in B,\\
  0&x'\notin B
 \end{cases}
\]
は $H^{1/2}(\mathbb R^{n-1})$ に属し，
$\norm{T_\eta g}_{H^{1/2}(\mathbb R^{n-1})}
 \le C\norm g_{H^{1/2}(\partial\Omega)}$ を満たす．
逆に $\operatorname{supp}h\subset K$ である
$h\in H^{1/2}(\mathbb R^{n-1})$ を
$V$ 上で $h\circ\psi^{-1}$，$V$ の外で零として
得られる $\widetilde h$ は
$H^{1/2}(\partial\Omega)$ に属し，
\[
 \norm{\widetilde h}_{H^{1/2}(\partial\Omega)}
 \le C\norm h_{H^{1/2}(\mathbb R^{n-1})}
\]
を満たす．
\end{lemma}
\begin{proof}
$m=n-1$ と置く．有界Lipschitz境界は
有限個のLipschitzグラフで覆われるから，
ある $r_0,C>0$ に対して
\[
 dS(\partial\Omega\cap B(y,r))\le Cr^m
 \quad(y\in\partial\Omega,\ 0<r<r_0)
\]
が成立する．実際，各グラフ上の面積要素は
$dx'$ の有界倍であり，グラフ上の距離が $r$ 以下なら
射影先の距離も $r$ 以下である．有限被覆を用いれば
同じ定数で評価できる．
従って $2^{-k-1}r_0\le|x-y|<2^{-k}r_0$
という環状領域で積分を見積もると
\[
 \sup_{y\in\partial\Omega}
 \int_{\partial\Omega}
 \frac{\min(1,|x-y|^2)}{|x-y|^{m+1}}\,dS_x
 <\infty.
 \tag{$*$}
\]
実際，原点近くの第 $k$ 環の積分は
$C2^{-k}r_0$ 以下であり，
$r_0$ より遠い部分は境界面積が有限であるため
積分可能である．

$\eta$ は有界かつLipschitzであるから
\[
 |\eta(x)g(x)-\eta(y)g(y)|^2
 \le 2\norm\eta_\infty^2|g(x)-g(y)|^2
   +C|g(y)|^2\min(1,|x-y|^2).
\]
これを境界上で積分し，$(*)$ を適用すると
乗算の有界性を得る．

次に $\psi$ はbi-Lipschitzであり，
面積公式から
\[
 dS_{\psi(x')}=
 J_\varphi(x')\,dx',\qquad
 1\le J_\varphi(x')
 =\sqrt{1+|\nabla\varphi(x')|^2}\le C.
\]
従って $B\times B$ 上の差分積分は
$V\times V$ 上の差分積分と
上下の定数倍を除き一致する．
また $\operatorname{supp}(T_\eta g)\subset K$ であり，
$d=\operatorname{dist}(K,\mathbb R^m\setminus B)>0$
だから，
\[
 \sup_{x'\in K}\int_{\mathbb R^m\setminus B}
       |x'-y'|^{-m-1}\,dy'<\infty.
\]
ゆえに $B$ の内外を結ぶ差分積分は
$C\norm{T_\eta g}_{L^2(B)}^2$ 以下となる．
境界上での乗算評価と合わせて
$T_\eta g$ のノルム評価を得る．

逆向きにも $V\times V$ の差分積分は
$B\times B$ のものと比較できる．
$\psi(K)$ は境界の相対開集合 $V$ の内部に
コンパクトに含まれるから
$d'=\operatorname{dist}(\psi(K),
 \partial\Omega\setminus V)>0$
（$V=\partial\Omega$ なら交差積分は不要）である．
従って $V$ と $\partial\Omega\setminus V$ の
交差積分は $C\norm h_{L^2}^2$ 以下である．
これで逆向きの評価も従う．
\end{proof}

\section{Lipschitz領域の一般トレース定理}
\begin{theorem}[Sobolevトレース定理]
$\Omega\subset\mathbb R^n$ を有界Lipschitz領域とする．
次を満たす有界線形作用素が存在する．
\[
 \operatorname{Tr}:H^1(\Omega)\to H^{1/2}(\partial\Omega),
 \qquad\operatorname{Tr}u=u|_{\partial\Omega}
 \quad(u\in C^\infty(\overline\Omega)).
\]
この作用素は全射であり，有界線形右逆
$R:H^{1/2}(\partial\Omega)\to H^1(\Omega)$ を持つ．
\end{theorem}
\begin{proof}
$\partial\Omega$ の有限個のグラフ座標を
$Q_j=\Phi_j(B_j\times(-h_j,h_j))$ とし，
$\psi_j(x')=\Phi_j(x',0)$ と置く．
$\operatorname{supp}\eta_j\Subset Q_j$ となる
滑らかな非負函数 $\eta_j$ を選び，
$\sum_j\eta_j=1$ が $\partial\Omega$ 上で
成立するようにする．各 $\eta_j|_{\partial\Omega}$ の台は
$\psi_j(K_j)$（$K_j\Subset B_j$）に含まれるように取る．

まず $u$ が $\overline\Omega$ の近傍で
滑らかな場合を考える．
\[
 v_j(x',s)=
 \bigl(\eta_ju\bigr)(\Phi_j(x',s)),
 \quad (x',s)\in B_j\times(0,h_j)
\]
と置き，半空間の残りでは零とする．
補題\ref{lem:chart-H1}によって
$v_j\in H^1(\mathbb R^n_+)$ で
$\norm{v_j}_{H^1}\le C_j\norm u_{H^1(\Omega)}$
である．
$u$ と $\eta_j$ は境界まで連続なので，
\[
 v_j(\cdot,s)\longrightarrow
 ((\eta_j u)\circ\psi_j)^0
 \quad\text{in }L^2(\mathbb R^{n-1})
 \qquad(s\downarrow0)
\]
が成り立つ．ここで上添字 $0$ は
$B_j$ の外での零延長を意味する．
実際，$\eta_j$ の台はチャートの内部に
コンパクトに含まれるので，断面上での収束は
共通のコンパクト集合上の一様収束から従う．
補題\ref{lem:halfspace-slice}により
\[
 \operatorname{Tr}_+v_j=((\eta_j u)\circ\psi_j)^0.
 \tag{$*$}
\]
半空間のトレース不等式と
補題\ref{lem:boundary-local}の逆向きの評価から
\[
 \norm{\eta_j u|_{\partial\Omega}}_{H^{1/2}(\partial\Omega)}
 \le C_j\norm{\operatorname{Tr}_+v_j}_{H^{1/2}(\mathbb R^{n-1})}
 \le C_j'\norm u_{H^1(\Omega)}
\]
を得る．有限和 $u|_{\partial\Omega}
=\sum_j\eta_j u|_{\partial\Omega}$ を用いると
\[
 \norm{u|_{\partial\Omega}}_{H^{1/2}(\partial\Omega)}
 \le C_\Omega\norm u_{H^1(\Omega)}.
\]
補題\ref{lem:lip-density}による稠密性から，
境界への制限は一意な有界線形作用素
$\operatorname{Tr}$ に延長される．

この延長の局所座標との両立性を確認する．
任意の $u\in H^1(\Omega)$ と
$\zeta\in C_c^\infty(Q_j)$ を取り，
$v=((\zeta u)\circ\Phi_j)^0$ とする．
補題\ref{lem:lip-density}により
滑らかな $u_k\to u$ in $H^1(\Omega)$ を取る．
局所座標変換の評価から
$v_k=((\zeta u_k)\circ\Phi_j)^0\to v$
in $H^1(\mathbb R^n_+)$ であり，
$(*)$ と同じ境界断面の議論から
\[
 \operatorname{Tr}_+v_k
 =((\zeta u_k)\circ\psi_j)^0.
\]
一方，$\operatorname{Tr}u_k\to\operatorname{Tr}u$
in $H^{1/2}(\partial\Omega)$ であり，
補題\ref{lem:boundary-local}により
$\zeta\operatorname{Tr}u_k$ の平坦座標への
零延長も $H^{1/2}$ で収束する．
よって両辺で極限を取ると
\[
 \operatorname{Tr}_+
  \bigl(((\zeta u)\circ\Phi_j)^0\bigr)
 =((\zeta\operatorname{Tr}u)\circ\psi_j)^0
 \tag{$**$}
\]
を得る．

次に全射性を示す．
$g\in H^{1/2}(\partial\Omega)$ とし，
\[
 g_j=((\eta_jg)\circ\psi_j)^0
 \in H^{1/2}(\mathbb R^{n-1})
\]
と置く．補題\ref{lem:boundary-local}により
$\norm{g_j}_{H^{1/2}}\le C_j\norm g_{H^{1/2}(\partial\Omega)}$
である．半空間の右逆を用いて
$w_j=Eg_j\in H^1(\mathbb R^n_+)$ とする．
$\chi_j\in C_c^\infty(Q_j)$ を
$\chi_j=1$ が
$\operatorname{supp}(\eta_j|_{\partial\Omega})$
の近傍で成り立つように取る．
局所座標変換の補題によって
\[
 u_j(x)=
 \begin{cases}
 \chi_j(x)w_j(\Phi_j^{-1}(x))
       &x\in\Omega\cap Q_j,\\
 0     &x\in\Omega\setminus Q_j
 \end{cases}
\]
は $H^1(\Omega)$ に属し，
\[
 \norm{u_j}_{H^1(\Omega)}
 \le C_j\norm{w_j}_{H^1(\mathbb R^n_+)}
 \le C_j'\norm g_{H^{1/2}(\partial\Omega)}
\]
を満たす．

そのトレースを同定する．
$w_j$ の $L^2$ 値断面は
補題\ref{lem:halfspace-slice}から
$w_j(\cdot,s)\to g_j$ in $L^2$ である．
$\zeta\in C_c^\infty(Q_j)$ とすると，
$\bigl(\zeta\chi_j\bigr)\circ\Phi_j$ は
$s\downarrow0$ で
$\bigl(\zeta\chi_j\bigr)\circ\psi_j$ に一様収束し，
全て一様有界なので
\[
 \operatorname{Tr}_+
  \bigl(((\zeta u_j)\circ\Phi_j)^0\bigr)
 =\bigl((\zeta\chi_j)\circ\psi_j\bigr)g_j.
\]
これと $(**)$ を比較すると，
$\operatorname{Tr}u_j=\chi_j\,\eta_jg
=\eta_jg$ が $\partial\Omega\cap Q_j$ で成り立つ．
また $u_j$ は $Q_j$ の外の近傍では零なので，
同じ局所恒等式を他のチャートに適用すれば
$\operatorname{Tr}u_j=0$ が
$\partial\Omega\setminus Q_j$ で成り立つ．
従って全境界上で
$\operatorname{Tr}u_j=\eta_jg$ である．

最後に $Rg=\sum_j u_j$ と定めれば，
有限和なので $R$ は線形で，
\[
 \norm{Rg}_{H^1(\Omega)}
 \le C\norm g_{H^{1/2}(\partial\Omega)},\qquad
 \operatorname{Tr}(Rg)=\sum_j\eta_jg=g
\]
を得る．よって $R$ は有界線形右逆であり，
$\operatorname{Tr}$ は全射である．
\end{proof}

\section{零トレースと零境界Sobolev空間}
\begin{lemma}[Lipschitz領域上のGauss--Green公式]
\label{lem:lip-green}
$u$ が $\overline\Omega$ の近傍で滑らか，
$\phi\in C_c^\infty(\mathbb R^n)$ とする．
$\nu$ を $\partial\Omega$ の外向き単位法線
（ほとんど至る所で定義される）とすると，
各 $j=1,\ldots,n$ について
\[
 \int_\Omega u\,\partial_j\phi\,dx
 =-\int_\Omega(\partial_ju)\phi\,dx
 +\int_{\partial\Omega}u\phi\,\nu_j\,dS
\]
が成り立つ．
\end{lemma}
\begin{proof}
積の微分公式により，
$\int_\Omega\partial_j(u\phi)
=\int_{\partial\Omega}u\phi\,\nu_j\,dS$
を示せばよい．
境界の有限個のグラフ座標と内部の開集合に
従属する滑らかな分割単位を用いると，
$F=u\phi$ を有限個の局所成分に分解できる．
内部に台を持つ成分の導函数の積分は零である．
境界の一つの成分を再び $F$ と書き，
その台が座標の側面と上面から離れているものとする．
剛体運動後，$\Omega$ は
$t>\varphi(x')$ で表され，
$\varphi$ はLipschitzである．

まず法線方向には，一変数の積分の基本定理から
\[
 \int_{\Omega}\partial_tF\,dx'\,dt
 =-\int_B F(x',\varphi(x'))\,dx'
\]
を得る．
接線座標 $x_i$（$1\le i<n$）については
\[
 G(x')=\int_{\varphi(x')}^\infty F(x',t)\,dt
\]
を考える．$G$ はコンパクト台を持つ
Lipschitz函数であり，
Rademacherの定理と積分の微分法則によって
ほとんど至る所で
\[
 \partial_iG(x')
 =\int_{\varphi(x')}^\infty\partial_iF(x',t)\,dt
 -F(x',\varphi(x'))\partial_i\varphi(x')
\]
が成り立つ．$G$ の弱導函数を全空間上で
積分すると零だから
\[
 \int_\Omega\partial_iF\,dx
 =\int_BF(x',\varphi(x'))\partial_i\varphi(x')\,dx'.
\]
グラフの外向き単位法線は
\[
 \nu=\frac{(\nabla\varphi,-1)}
           {\sqrt{1+|\nabla\varphi|^2}},
 \qquad
 dS=\sqrt{1+|\nabla\varphi|^2}\,dx'
\]
であるから，以上の二式はそれぞれ
$\int_{\partial\Omega}F\nu_j\,dS$ に等しい．
剛体運動前の座標に戻して有限個の成分を加えれば，
全ての $j$ について公式が得られる．
\end{proof}

\begin{lemma}[零延長]
$u\in H^1(\Omega)$ かつ $\operatorname{Tr}u=0$ とする．
$u$ を $\mathbb R^n\setminus\Omega$ で零とした函数
$\widetilde u$ は $H^1(\mathbb R^n)$ に属する．
\end{lemma}
\begin{proof}
稠密性の補題から $u_k\in C^\infty(\overline\Omega)$ で
$u_k\to u$ in $H^1(\Omega)$ となる列を選ぶ．
境界トレースの連続性から
$u_k|_{\partial\Omega}\to0$ in $L^2(\partial\Omega)$ である．
$\varphi\in C_c^\infty(\mathbb R^n)$ に対し，
補題\ref{lem:lip-green}を $u_k$ に適用すると
\[
 \int_\Omega u_k\partial_j\varphi
 =-\int_\Omega (\partial_j u_k)\varphi
   +\int_{\partial\Omega}u_k\varphi\,\nu_j\,dS
\]
が成立する．
$k\to\infty$ で境界積分はCauchy--Schwarzによって0となる．
したがって
\[
 \int_{\mathbb R^n}\widetilde u\partial_j\varphi
 =-\int_{\mathbb R^n}
   \widetilde{\partial_j u}\,\varphi
\]
が成り立ち，零延長の弱微分は弱微分の零延長に等しい．
双方 $L^2$ に属するから $\widetilde u\in H^1(\mathbb R^n)$ である．
\end{proof}

\begin{theorem}[トレースの核]
有界Lipschitz領域 $\Omega$ について
\[
 \ker\operatorname{Tr}=H_0^1(\Omega)
\]
が成り立つ．
\end{theorem}
\begin{proof}
$H_0^1(\Omega)$ は $C_c^\infty(\Omega)$ の
$H^1$ 閉包であり，
$C_c^\infty(\Omega)$ の函数は境界近傍で零である．
トレースの連続性から
$H_0^1(\Omega)\subset\ker\operatorname{Tr}$ である．

逆向きに $\operatorname{Tr}u=0$ とし，
前補題によって得られた
零延長 $\widetilde u\in H^1(\mathbb R^n)$ を取る．
$\overline\Omega$ を覆う内部の開集合と有限個の
境界グラフ座標近傍に従属する滑らかな分割単位を取り，
$\widetilde u=\sum_{j=0}^N v_j$，
$v_j=\eta_j\widetilde u$ とする．
各 $v_j$ は $H^1(\mathbb R^n)$ に属し，
$\operatorname{supp}v_j$ は対応する近傍の内部に
コンパクトに含まれる．
内部成分 $v_0$ は $\Omega$ の内部に台を持つから，
通常のmollifierとの畳み込みで
$C_c^\infty(\Omega)$ 函数に近似できる．

境界近傍の一つを考え，
剛体運動後の座標を $(x',t)$ とする．
そこでは $\Omega=\{t>\varphi(x')\}$ であり，
$v_j$ は $t\le\varphi(x')$ で
ほとんど至る所零である．
$\varepsilon>0$ に対して，全空間上で
\[
 v_j^\varepsilon(x',t)
  =v_j(x',t-\varepsilon)
\]
と定める．全空間 $H^1$ の平行移動連続性から
\[
 \norm{v_j^\varepsilon-v_j}_{H^1(\mathbb R^n)}
 \longrightarrow0\qquad(\varepsilon\downarrow0).
\]
また $v_j^\varepsilon$ の本質的な台は
局所座標内で
$t\ge\varphi(x')+\varepsilon$ に含まれる．
$\varphi$ のLipschitz定数を $L$ とし，
\[
 0<\delta<\frac{\varepsilon}{2(1+L)}
\]
を満たす $\delta$ のmollifierを取る．
$\rho_\delta*v_j^\varepsilon$ が点 $(x',t)$ で
零でないなら，ある $(y',s)$，
$|(y',s)|<\delta$ に対して
\[
 t-s\ge\varphi(x'-y')+\varepsilon
 \ge\varphi(x')-L\delta+\varepsilon
\]
である．従って
\[
 t-\varphi(x')\ge\varepsilon-(1+L)\delta
 >\varepsilon/2.
\]
この評価により畳み込みの台は
境界から正の距離を持って $\Omega$ 内に含まれる．
$\varepsilon,\delta$ をチャートの台の余裕よりも
十分小さく取れば，側面・上面に達することもない．
従って
$\rho_\delta*v_j^\varepsilon\in C_c^\infty(\Omega)$
であり，
\[
 \norm{\rho_\delta*v_j^\varepsilon-v_j}_{H^1(\mathbb R^n)}
 \le
 \norm{\rho_\delta*v_j^\varepsilon-v_j^\varepsilon}_{H^1(\mathbb R^n)}
 +\norm{v_j^\varepsilon-v_j}_{H^1(\mathbb R^n)}
\]
となる．最初に $\varepsilon$，次に $\delta$ を
十分小さくすればこの右辺は任意に小さくなる．

有限個の局所成分をそれぞれ近似して加えると，
$C_c^\infty(\Omega)$ 函数列で $\widetilde u$ を
全空間 $H^1$ 近似できる．$\Omega$ に制限すれば
$u\in H_0^1(\Omega)$ を得る．
\end{proof}

\section{零境界条件の解釈}
\begin{corollary}
有界Lipschitz領域 $\Omega$ 上で $u\in H^1(\Omega)$ に対し，
$u\in H_0^1(\Omega)$ であることと
$\operatorname{Tr}u=0$ であることは同値である．
特に $u$ が $\overline\Omega$ 上で連続な代表を持つ場合，
$u\in H_0^1(\Omega)$ であることは，その連続代表の
境界値が $\partial\Omega$ の各点で零であることと同値である．
\end{corollary}
\begin{proof}
前定理から第一の同値性は直ちに従う．
$u$ が境界まで連続なら，グラフ座標での
$u(x',\varphi(x')+t)$ の $t\downarrow0$ における
$L^2$ 極限は境界値に等しい．一方，一次元の端点値評価と
局所化によってこの極限は $\operatorname{Tr}u$ に等しい．
従って境界値が $L^2(\partial\Omega)$ で零なら，
連続性と各開境界部分の正の面積から境界上の全点で零となる．
\end{proof}

\begin{remark}
本稿は $H^1$ の境界トレースの完全な構成を与える．
高階の $H^s$ トレースの一般形，角点における
高階微分の境界条件，任意Lipschitz領域での楕円方程式の
$H^2$ 正則性は別の問題であり，本稿からは従わない．
\end{remark}

\begin{thebibliography}{9}
\bibitem{adams} R. A. Adams and J. J. F. Fournier,
 \emph{Sobolev Spaces}, 2nd ed., Academic Press, 2003.
\bibitem{leoni} G. Leoni, \emph{A First Course in Sobolev Spaces},
 2nd ed., AMS, 2017.
\bibitem{mclean} W. McLean,
 \emph{Strongly Elliptic Systems and Boundary Integral Equations},
 Cambridge University Press, 2000.
\end{thebibliography}

\section{到達課題の完全解答}
\noindent\textbf{一次元のトレースの計算．}
\begin{proposition}
$u(x)=x(1-x)$ は $H_0^1(0,1)$ に属し，
端点のトレースはともに零である．
一方 $v(x)=x$ は $H^1(0,1)$ に属するが
$H_0^1(0,1)$ に属さない．
\end{proposition}
\begin{proof}
両者は閉区間上滑らかなので $H^1(0,1)$ に属する．
まず $u$ が $H_0^1(0,1)$ に属することを
直接示す．$\vartheta\in C^\infty(\mathbb R)$ を
$0\le\vartheta\le1$，
$\vartheta(t)=0$（$t\le1$），
$\vartheta(t)=1$（$t\ge2$）
となるように取り，
\[
 \chi_\varepsilon(x)
 =\vartheta(x/\varepsilon)
  \vartheta((1-x)/\varepsilon),
 \qquad0<\varepsilon<1/4
\]
と置く．
$\chi_\varepsilon u\in C_c^\infty(0,1)$ である．
$u(0)=u(1)=0$ と $u$ の滑らかさから
\[
 |u(x)|\le C\min(x,1-x)
\]
が成立する．
$(1-\chi_\varepsilon)u$ とその導函数は
$(0,2\varepsilon)\cup(1-2\varepsilon,1)$
の外で零であり，
$|\chi_\varepsilon'|\le C/\varepsilon$ なので
\[
 \norm{(1-\chi_\varepsilon)u}_{L^2}
 \le C\varepsilon^{3/2},\qquad
 \norm{\bigl((1-\chi_\varepsilon)u\bigr)'}_{L^2}
 \le C\varepsilon^{1/2}.
\]
よって $\chi_\varepsilon u\to u$ in $H^1(0,1)$
となり，$u\in H_0^1(0,1)$ である．

他方，一次元Sobolev函数は絶対連続な代表を持ち，
微積分学の基本定理から任意の
$h\in H^1(0,1)$ に対して
\[
 |h(1)|\le\norm h_{L^2(0,1)}
                 +\norm{h'}_{L^2(0,1)}
\]
が成り立つ．実際
$h(1)=h(t)+\int_t^1h'(s)\,ds$ を
$t\in(0,1)$ で積分し，絶対値を評価すればよい．
従って端点評価 $h\mapsto h(1)$ は
$H^1$ 上の連続線形汎函数であり，
$C_c^\infty(0,1)$ の全ての元を零へ写す．
その閉包 $H_0^1(0,1)$ の全ての元も
$h(1)=0$ を満たすが，$v(1)=1$ である．
従って $v\notin H_0^1(0,1)$ である．
\end{proof}
\noindent\textbf{到達確認．}
$w(x)=\sin(\pi x)$ は $H_0^1(0,1)$ に属するか．
\par\noindent\textbf{解答．}
$w$ は閉区間上滑らかで $w(0)=w(1)=0$ だから，
上の証明と同じ切断函数 $\chi_\varepsilon$ を用いれば
$\chi_\varepsilon w\to w$ in $H^1(0,1)$ を得る．
従って $w\in H_0^1(0,1)$ である．
\end{document}